% LNCS submission — Springer Computer Science Proceedings.
% Paper A of a two-paper plan: the meta-science system paper. The companion
% instrument paper ("Recall or Discovery?") lives on branch paper-instrument.
% Rule of the pair: no result is reported in both.
%
% Style per the author guidelines (Hofmann et al.): abstract 70-150 words;
% middot keywords; two numbered heading levels; table captions above; numbered
% equations without section counters; [n] citations; Springer references;
% American English.
\documentclass{llncs}

\usepackage[T1]{fontenc}
\usepackage{amsmath}
\usepackage{amssymb}
\usepackage{booktabs}
\usepackage{graphicx}
\usepackage{url}

\begin{document}

\title{Meta-Science: Falsifiable Scientific Reasoning by Agents\\
on Anonymized Causal Worlds}

\author{Marco Vanadia}
\institute{Independent Researcher, Warsaw, Poland\\
\email{marco.vanadia@gmail.com}}

\maketitle

\begin{abstract}
``Self-improving agent'' is the most repeated unfalsifiable claim in the
field: a system that reports its own improvement is indistinguishable from one
that logs ``improved'' and changes nothing. We present a system built around
the opposite artifact---one that refuses. An agent does science on randomly
generated causal mini-worlds it has never seen: it commits falsifiable
predictions before each experiment, is refuted by its own interventions, and
proposes changes to its own method that are promoted only when they beat the
incumbent by a margin on held-out worlds the proposer cannot enumerate.
Anonymization is a tested property, not an assertion: reading observations
alone, a frontier model recovers the causal direction in zero of four
confounded worlds; allowed to intervene, four of four. Across three published
live runs the gate promoted once and refused twice---every time---and each
verdict carries a receipt sufficient to recompute it without trusting us.
\keywords{Falsifiability \and Agents \and Causal inference \and
Self-improvement \and Benchmarks \and Scientific method}
\end{abstract}


\section{Introduction}

Karl Popper's demarcation criterion applies to demonstrations of artificial
scientists as much as to theories~\cite{popper}: a demonstration that cannot
fail is empty. Yet demonstrations of ``self-improving'' agents typically
exhibit successes on benchmarks chosen after the fact, judged by the agent
itself or by nothing at all. The interesting artifact, we argue, is not a
system that improves; it is a system that \emph{refuses to}---and can show,
number by number, why.

This paper describes such a system, built on two refusals stated up front:
\begin{equation}
A \vdash c \;\;\not\Rightarrow\;\; \vDash c ,
\end{equation}
\emph{no agent is the judge of its own claims}---that an agent derives a
conclusion never entails that the conclusion holds, so verdicts are computed,
never solicited; and
\begin{equation}
P(Y \mid \mathrm{do}(X)) \;\not\equiv\; P(Y \mid X) ,
\end{equation}
\emph{seeing is not doing}~\cite{pearl}---the two coincide exactly when
confounding is absent, and that absence may never be assumed. The first axiom
becomes an architecture (Sect.~\ref{sec:gate}); the second becomes a world
generator whose observational surface actively lies (Sect.~\ref{sec:worlds}).

The system operates on two levels under one gate. Level~1: an agent meets a
world it has never seen---opaque variable names, no documentation, two
affordances (look, or act)---forms a falsifiable hypothesis, designs an
intervention, and finds out whether it was wrong. Level~2: the agent proposes
a change to its own experiment-design strategy, evaluated as champion versus
challenger on held-out worlds it never sees, promoted only past a margin, with
an independent auditor free to dissent in writing and every verdict written to
a replayable receipt.

Our contributions: (i)~a generator of anonymized causal worlds whose
anonymization is \emph{tested} rather than asserted
(Sect.~\ref{sec:worlds}); (ii)~a promotion gate for self-modification with
margins against noise-ratcheting, receipts for lineage, and an advisory
auditor against metric gaming (Sect.~\ref{sec:gate}); (iii)~measurements,
including a 384-run study of discovery by refutation and three published live
runs of gated self-evolution whose headline behaviors reproduce run over run
(Sect.~\ref{sec:results}); and (iv)~a methodological confession log---our
benchmark flattered us once, our published figures drifted from the code once,
and both failures, with their fixes, are part of the record
(Sect.~\ref{sec:honesty}).

\section{Related Work}

Automated scientific discovery has substantial prior art: active
learning~\cite{settles}, constraint-based causal discovery~\cite{spirtes},
distilling laws from data~\cite{schmidt}, and rule induction in the
Zendo/Eleusis tradition of hidden-rule games~\cite{gardner}. Recent
LLM-as-scientist systems~\cite{aiscientist} close the loop from hypothesis to
paper but are evaluated on their successes. This is not the first artificial
scientist and does not claim to be. The narrower claim is \emph{falsifiable
self-improvement}: a loop whose every promotion is earned against evidence the
proposer cannot see, score, or tune against---leaving a receipt either way. A
companion paper~\cite{companion} builds a recall-versus-inference instrument
on this framework; no result is reported in both.

\section{Anonymized Causal Worlds}\label{sec:worlds}

Worlds are structural causal models seeded from real scientific
topologies---state relations, exponential response, transport chains,
compartment flows, inheritance, confounding---because real science supplies
non-trivial structure and known ground truth for free. Six frozen templates
rotate by seed; every retrievable surface is stripped at birth: domain names
become $X_1..X_n$, labels shuffle per world, constants and exponents
randomize, and the functional family varies. Generation is deterministic from
the seed via a stable hash, so the same seed builds a bit-identical world in
every process---asserted in a subprocess, because an in-process check cannot
see the failure mode (Sect.~\ref{sec:honesty}).

One template family is required to \emph{invert}: a hidden confounder gives
the observational correlation the opposite sign of the true causal effect
(Simpson-style), with inversion realized in 83.9\,\% of generated confounded
worlds across the study. On such worlds an agent that only looks concludes
the opposite of the truth.

\textbf{Anonymization as a tested property.} A banned lexicon is necessary
but not sufficient---the real question is whether a model can answer without
experimenting. So it is tested adversarially: give the model the anonymized
observations, forbid experiments, and ask for the causal direction
(Table~\ref{tab:blind}). Observation alone recovers zero of four confounded
worlds; the same loop allowed to intervene recovers four of four. The
comparison shows simultaneously that the disguise holds (if the model could
recognize the worlds, looking would have sufficed) and that intervention does
real work rather than decorating a conclusion the model already had.

\begin{table}
\caption{The measurement that anchors both axioms: identical anonymized
worlds; the only difference is the ability to act.}\label{tab:blind}
\centering
\begin{tabular}{@{}ll@{}}
\toprule
 & Confounded worlds recovered\\
\midrule
Frontier model, observation only & 0 / 4\\
The same loop, allowed to intervene & 4 / 4\\
\bottomrule
\end{tabular}
\end{table}

\section{The Gate}\label{sec:gate}

Every record lives in one of three tiers---raw (anything the agent produces),
wiki (candidate findings), output (accepted canon)---and a record's tier
\emph{is} its status: nothing is promoted by assertion, only by moving up.
The agent never sees the held-out seeds, never runs the scorer, and never
writes above raw. Promotion of a challenger method $M'$ over champion $M$
requires
\begin{equation}
M' \succ M \iff S(M', W) \ge S(M, W) + \varepsilon,
\qquad W \cap \mathrm{view}(M') = \varnothing ,
\end{equation}
with $\varepsilon = 0.02$ on a score that separates accuracy from measurement
cost. The margin exists because a gate without one ratchets on noise: a
thousand tiny lucky wins launder randomness into ``progress.'' Every verdict,
promotion and refusal alike, writes a receipt carrying the diff, both scores,
the margin, the score decomposition, and the held-out seeds---enough to
recompute the decision independently, and cheap enough that refusals are
documented as thoroughly as successes.

An \emph{independent auditor}---a different model---reads each promotion's
receipt and may dissent on the record, without veto. Advisory status is
deliberate: the auditor is a language model and not deterministic
(Sect.~\ref{sec:results}); a gate whose verdicts depended on it would inherit
that variance.

\section{Measurements}\label{sec:results}

\textbf{Discovery by refutation.} Predictions are committed to the trace
before each experiment runs; the verdict compares prediction to result, and
the model is never asked whether it was right---a hypothesis recorded after
the result is a description, not a hypothesis. Across a 384-run study the
loop attains recall 1.000 and precision 0.786 on true causal edges, and an
observational pre-screening heuristic shows no benefit over direct
experimentation (a null result, reported as such).

\textbf{Gated self-evolution, three published runs.} From the same starting
champion and note, three consecutive live runs were recorded and all three
published---none discarded, none re-rolled (Table~\ref{tab:runs}). The shape
reproduces exactly: one promotion, two refusals, every time. The promoted
change was always a large cut in measurement (samples per arm 400 to 150--200);
everything after it was refused---four of the six refusals were of candidates
that genuinely scored \emph{higher} and missed the margin, which is the gate
doing the one thing it exists to do. In every run the proposer, given
structured verdict history, stopped repeating refused proposals and changed
knobs by generation three.

\begin{table}
\caption{Three live runs of gated self-evolution. Scores are on 24 held-out
worlds; margin $\varepsilon=0.02$. The auditor disagreement in runs 1 and 3
(identical promotion, opposite audit) is retained, not resolved.}\label{tab:runs}
\centering
\begin{tabular}{@{}llll@{}}
\toprule
 & Run 1 & Run 2 & Run 3\\
\midrule
Gen 1 & promoted $+0.0416$ & promoted $+0.0277$ & promoted $+0.0416$\\
Gen 2 & refused $+0.0084$ & refused $+0.0139$ & refused $+0.0139$\\
Gen 3 & refused $-0.0287$ & refused $-0.0259$ & refused $+0.0042$\\
Audit & flagged & flagged & legitimate\\
\bottomrule
\end{tabular}
\end{table}

\textbf{What does not reproduce is reported.} Runs 1 and 3 contain the
identical promotion---same diff, same scores---and the auditor flagged it in
one and called it legitimate in the other. The disagreement is preserved on
both receipts. The margin decides either way; the auditor annotates the
record rather than controlling it.

\section{The Process Was the Philosophy}\label{sec:honesty}

Building the system forced its discipline onto its builders, and the failures
are part of the evidence.

\textbf{Our benchmark flattered us.} Intervention arms initially shared
random draws (common random numbers), which cancels noise so completely that
cutting measurement was free score---very likely why the proposer kept
winning by asking for fewer samples. We measured both regimes: at 25 samples
per arm, accuracy holds at 98.1\,\% under shared noise and falls to 84.3\,\%
under honest independent noise. Independent noise is now the default, the
extreme strategy correctly loses, and both facts are asserted by tests.

\textbf{Our published figures drifted from the code.} World generation was
seeded with a per-process randomized hash, so the same seed built a different
world in every process and the replayability the receipts promise did not
hold. The fix (a stable content hash) is guarded by a subprocess determinism
test, and every published figure is pinned to the code by a test that fails
when either moves.

\textbf{Our first metric was gameable.} Cost initially charged the number of
experiments, leaving sample size free; a challenger bought accuracy with data
and appeared to have learned. Cost now charges total measurement, and that
candidate loses.

None of this was planned; all of it is the point. A method applied only to
one's subject is a pose. Applied to oneself, it becomes a practice.

\section{Limitations}

The live evidence is three runs of three generations with one model family
proposing; variance estimates need a repetition campaign. The worlds'
mechanisms are simple (the study's replication analysis found repetition
cost-neutral under i.i.d.\ Gaussian noise---worlds kind enough that
replication has nothing to catch, which is a statement about the worlds, not
about replication). The auditor's non-determinism is documented rather than
solved. The system recovers causal directions and magnitudes, not yet laws;
law induction is the companion paper's subject~\cite{companion}.

\section{Conclusion}

We set out to build not a self-improving agent but the gate such an agent
would have to pass, and only then the agent. The result improves rarely and
refuses often, which is precisely its value: the refusals carry information,
the promotions carry receipts, and every number in this paper can be
recomputed from committed artifacts without trusting its authors.
\emph{Ex probatione propria non sequitur veritas}: from one's own proving,
truth does not follow. The gate exists so that it does not have to.

\subsubsection*{Acknowledgements.} The system was built meta-agentically: an
AI pair-engineer proposing at speed, a human judging, refusing, and
redirecting---neither trusted as the judge of its own claims.

\begin{thebibliography}{9}
% NOTE: verify every entry against the actual editions before submission.
\bibitem{popper}
Popper, K.: Logik der Forschung. Julius Springer, Vienna (1934). English
edition: The Logic of Scientific Discovery. Hutchinson, London (1959)

\bibitem{pearl}
Pearl, J.: Causality: Models, Reasoning, and Inference. 2nd edn. Cambridge
University Press, Cambridge (2009)

\bibitem{settles}
Settles, B.: Active Learning. Synthesis Lectures on Artificial Intelligence
and Machine Learning. Morgan \& Claypool (2012)

\bibitem{spirtes}
Spirtes, P., Glymour, C., Scheines, R.: Causation, Prediction, and Search.
2nd edn. MIT Press, Cambridge (2000)

\bibitem{schmidt}
Schmidt, M., Lipson, H.: Distilling free-form natural laws from experimental
data. Science 324(5923), 81--85 (2009)

\bibitem{gardner}
Gardner, M.: Mathematical games: about the game of Eleusis, in which players
frame laws and solve them by induction. Scientific American 200(6), 160--168
(1959)

\bibitem{aiscientist}
Lu, C., Lu, C., Lange, R.T., Foerster, J., Clune, J., Ha, D.: The AI
Scientist: towards fully automated open-ended scientific discovery.
arXiv:2408.06292 (2024)

\bibitem{companion}
Vanadia, M.: Recall or Discovery? A Falsifiable Instrument for Law Induction
with Language Models. Companion paper, in preparation (2026)
\end{thebibliography}

\end{document}
